\documentclass[10pt,a4paper,sans]{moderncv}
\moderncvstyle{classic}
\moderncvcolor{blue}
\usepackage[utf8]{inputenc}
\usepackage[scale=0.83]{geometry} % scale=0.83 (vs 0.8) keeps this CV to two pages
\usepackage{helvet}

\name{Alexandro}{Disla}
\title{Officier de Suivi \& Évaluation | Systèmes d'Information Sanitaire \& Qualité des Données}
\address{Rue Saint-Louis Jeanty \#14, Route de Frère}{Port-au-Prince, Haïti}{}
\phone[mobile]{(+509) 4148-3700}
\email{alexandrodisla@hotmail.com}
\extrainfo{GitHub: \url{https://github.com/AD0791}}

\begin{document}
\makecvtitle

\section{Profil Professionnel}
\cvitem{}{Professionnel du suivi-évaluation avec plus de cinq ans d'expérience en Haïti, dont trois comme \textbf{Officier de Suivi \& Évaluation sur un programme VIH}, où j'ai conduit le rapportage de routine de cliniques multi-sites et rapporté les \textbf{indicateurs MER à l'USAID via DATIM (DHIS2)}. Mon travail porte sur la part du suivi-évaluation qui doit résister à la vérification : protocoles d'assurance qualité des données (DQA) établis et supervisés, chiffres remontés jusqu'au registre source, et collecte numérique standardisée afin que la donnée soit juste à la saisie plutôt que réparée ensuite. Économiste appliqué de formation et \textbf{ingénieur de données} en parallèle, j'administre également les bases et les systèmes de collecte dont dépendent les institutions sanitaires, et pas seulement les indicateurs qu'ils produisent. Résidant à Port-au-Prince ; français, anglais et créole haïtien courants.}

\section{Compétences Clés}
\cvitem{}{
\begin{itemize}
    \item \textbf{Systèmes de suivi-évaluation \& rapportage de routine :} Plans de suivi, cadres logiques, chaînes de résultats et tableaux de suivi des indicateurs ; cycles de rapportage mensuel, trimestriel, semestriel et annuel depuis des institutions sanitaires multi-sites ; rapportage bailleur des indicateurs via DATIM/MER ; planification et gestion d'enquêtes ; analyse qualitative et quantitative avec les désagrégations exigées.
    \item \textbf{Systèmes d'information sanitaire \& qualité des données :} Protocoles DQA établis et supervisés, avec vérification par remontée aux registres sources et suivi des actions correctives jusqu'à clôture ; taux de complétude publié avec chaque indicateur ; administration des systèmes de collecte MySQL, CommCare, mWater et ODK ; déploiement des versions sur les sites, contrôle d'accès par rôle, documentation des incidents et stockage sécurisé.
    \item \textbf{Renforcement de capacités \& coordination :} Formation du personnel des institutions sanitaires, des enquêteurs et des agents de collecte aux outils, à la qualité des données et à l'analyse ; encadrement technique du personnel de suivi-évaluation et rédaction de guides techniques ; coordination avec les partenaires gouvernementaux (MSPP, DINEPA) et traduction des constats techniques pour une direction non technique.
\end{itemize}}

\section{Expériences en Suivi-Évaluation Sanitaire et Systèmes d'Information}

\cventry{Oct 2021 -- Août 2024}{Officier de Suivi \& Évaluation}{Impact Youth Project, Caris Foundation International}{Haïti}{}{
\begin{itemize}
    \item \textbf{Rapportage sanitaire de routine :} Conduite du cycle de suivi d'un programme VIH sur des cliniques multi-sites, avec consolidation des données des sites dans les délais et rapportage des \textbf{indicateurs MER à l'USAID via DATIM}, la plateforme de rapportage du PEPFAR bâtie sur DHIS2.
    \item \textbf{Systèmes de collecte :} Mise en place et administration des systèmes de collecte numérique qui alimentaient ce rapportage (CommCare, HIVHaiti, MySQL), avec standardisation des formulaires et des flux de données entre sites.
    \item \textbf{Assurance qualité des données :} Établissement et supervision des protocoles DQA sur les données intégrées MySQL, avec remontée des chiffres rapportés jusqu'aux registres sources et suivi des actions correctives jusqu'à leur clôture, et non jusqu'à leur simple consignation.
    \item \textbf{Encadrement :} Supervision des agents de collecte sur l'ensemble des sites et encadrement technique du personnel de suivi-évaluation sur les procédures et les contrôles qualité.
    \item \textbf{Reporting \& tableaux de bord :} Automatisation du reporting d'activité en Python et SQL, réduisant le temps de traitement manuel de 40 \% ; maintenance des tableaux de bord programme sur AWS (Quarto) et Power BI pour le suivi hebdomadaire des indicateurs par la direction de programme.
\end{itemize}}

\cventry{Jan 2026 -- Mars 2026}{Data Analyst \& Consultant MEAL}{Anseye Pou Ayiti}{Haïti (Télétravail)}{}{
\begin{itemize}
    \item Définition de la stratégie de collecte et mise en œuvre d'indicateurs de performance pour des programmes de leadership multisectoriels.
    \item Codage des formulaires XLSForm pour \textbf{ODK} et KoboToolbox avec logiques de saut et contraintes de validation, et maintenance des scripts d'intégration Python sur AWS.
    \item Supervision de la qualité des données remontées du terrain, administration du système ODK/Admin Panel et formation des agents de terrain aux outils de collecte.
\end{itemize}}

\cventry{Oct 2023 -- Jan 2026}{Consultant mWater \& Analyste de Données}{HANWASH}{Télétravail (Upwork)}{}{
\begin{itemize}
    \item \textbf{Conception du suivi-évaluation :} Raffinement des plans MEAL et des tableaux de suivi des indicateurs alignés sur les normes JMP pour des programmes d'eau et d'assainissement.
    \item \textbf{Gestion des données :} Déploiement d'enquêtes géospatiales complexes sur mWater, administration de la base incluant le nettoyage et la fusion des enregistrements doublons, et résolution des problèmes de collecte terrain pour préserver l'intégrité des données.
    \item \textbf{Automatisation \& formation :} Création de tableaux de bord automatisés (mWater, Power BI) connectés aux API de collecte pour un reporting en temps réel ; production de guides techniques et formation des équipes à leur usage.
\end{itemize}}

\cventry{Nov 2015 -- Août 2023}{Économiste -- Direction de Planification Économique et Sociale (DPES)}{Ministère de la Planification et de la Coopération Externe (MPCE)}{Port-au-Prince}{}{Contribution aux rapports de suivi de l'exécution du Plan Triennal d'Investissement 2014--2016 ; affectation à l'unité statistique et modélisation, avec participation à la construction du « Modèle de Planification du Développement », outil de prévision, de simulation et d'analyse d'impact des politiques publiques ; suivi des indicateurs macroéconomiques avec le Ministère de l'Économie et des Finances.}

\section{Expériences en Administration de Bases de Données et Systèmes}

\cventry{Févr 2026 -- Présent}{Ingénieur Logiciel (Fullstack)}{Tekkod LLC}{Télétravail}{}{Direction d'un programme de modernisation d'infrastructure mobile, avec fiabilité hors ligne et sécurité des données pour un usage terrain ; construction de tableaux de bord Looker Studio avec recherche et filtrage pour la visibilité opérationnelle.}

\cventry{Mars 2021 -- Mai 2024}{Développeur Backend \& Ingénieur de Données}{Tekkod LLC}{Télétravail}{}{Administration de bases relationnelles en accès concurrent et résolution des défaillances qui rendent un système inutilisable sur le terrain --- une page de connexion ramenée de 36 secondes à moins d'une demi-seconde par l'élimination de boucles de requêtes N+1, et des index composites ayant supprimé les conflits de verrouillage ; mise en place de pipelines ETL (Apache Beam) migrant les données MongoDB vers BigQuery, et de services REST sécurisés (FastAPI, JWT) avec contrôle d'accès par rôle.}

\cventry{Nov 2015 -- Mars 2021}{Consultant Logiciel \& Spécialiste Intégration de Données}{CassionSoft (Projet UNOPS)}{Haïti}{}{Travaux d'intégration de données sur un projet des Nations Unies, avec une exposition directe aux procédures et aux standards documentaires onusiens ; implémentation de systèmes d'authentification RBAC (OAuth2/JWT) et validation de la logique d'intégration par des tests unitaires rigoureux.}

\section{Formation}
\cvitem{2010--2014}{Diplôme d'Études Supérieures en Économie Appliquée --- C.T.P.E.A (Centre de Techniques de Planification et d'Économie Appliquée), Port-au-Prince}
\cvitem{2009--2010}{Baccalauréat (Mention) --- Institution Saint-Louis de Gonzague}
\cvitem{Autoformation}{FlatIron School (Software Engineering), CodeWithMosh, Udemy}

\section{Compétences Techniques}
\cvitem{Info. sanitaire}{\textbf{DHIS2 / DATIM (rapportage des indicateurs MER)}, CommCare, HIVHaiti, mWater, protocoles DQA, tableaux de suivi des indicateurs}
\cvitem{Collecte mobile}{\textbf{ODK (Open Data Kit), KoboToolbox, XLSForm}, logiques de saut et contraintes de validation, conception et gestion d'enquêtes}
\cvitem{Analyse \& reporting}{\textbf{MS Excel (avancé), MS Word, MS Access, MS PowerPoint, GSuite}, \textbf{Power BI (expert)}, SQL, Python (Pandas), R (Quarto), Stata, Looker Studio, Tableau, modélisation statistique et échantillonnage}
\cvitem{Infrastructure}{MySQL, PostgreSQL, MongoDB, SQLite, Azure, AWS, Google Cloud (BigQuery), Docker, Git}
\cvitem{Langues}{\textbf{Français (maternel)}, \textbf{Créole haïtien (maternel)}, \textbf{Anglais (courant)}}

\end{document}
